\answer{ex:20:1}{$t_{\text{v}}=\tau_r\ln\frac{1-L}{1-H}=16.8$~h, $t_{\text{s}}=\tau_d\ln\frac HL=5.8$~h, período $22.5$~h; o oscilador é puxado para $24$~h pela oscilação diária dos limiares: captura de fase~\eqref{eq:pll}.}
\answer{ex:20:2}{(a)~$L=0.111$; com $L=0.092$ o sono duraria $9.6$~h. (b)~$22\%$ ($1/0.82=1.22$), e não $18\%$.}
\answer{ex:20:3}{$H=\frac{p-1}{p-1/q}=0.623$, $L=H/q=0.093$, onde $p=e^{16/\tau_r}$, $q=e^{8/\tau_d}$. Depois de acordar após $4$~h, a fase se desloca para sempre (adormece $7.2$~h mais cedo); com limiares oscilantes, ela volta em dois ou três dias.}
\answer{ex:20:4}{(a)~$r_\pm=\frac12\bigl(1\pm\sqrt{1-4k/w}\bigr)=0.947,\ 0.053$; $a_{\text{s}}=3.51$, $a_{\text{i}}=0.49$. (b)~Atividade $\approx0.33$~s, silêncio $\approx0.59$~s, período $\approx0.93$~s, fração de atividade $0.36$.}
\answer{ex:20:5}{$a_{\text{s}}/r_+<b<a_{\text{i}}/r_-$: $3.71<b<9.20$. O \nt{ACET} reduz $b$ para menos de $3.71$: a interseção passa ao ramo superior, e o estado ativo ($r\approx1$) fica estável.}
\answer{ex:20:6}{$4.9$, $6.8$, $8.8$, $5.5$, $8.9$~h para $D=24$, $32$, $40$, $48$, $64$: a duração é definida pela fase diária do adormecer, e não pela dívida.}
\answer{ex:20:7}{Após B, o erro em A é em média $0.51$ (de $0.19$ a $1.08$ em $20$ conjuntos; a resposta nula dá $1$). Intercalando, os dois erros ficam abaixo de $10^{-4}$.}
\answer{ex:20:8}{Questão aberta. O escalonamento uniforme preserva as razões dos pesos, mas só preserva as respostas do neurônio de limiar se o limiar for escalonado pelo mesmo fator; os replays intercalados mudam as razões. A invariância dos contatos grandes contradiz um escalonamento puramente uniforme.}
